% sec07: comparison with the holomorphic twist (Budzik--Kulp). Body only.
% Notation per NOTATION.md: K has no repository superscript; T^{HT,...} is
% used only inside this section, where the comparison is defined.
\section{Comparison with the holomorphic twist}\label{sec:ht}

The one-loop coefficients of Section~\ref{sec:superspace} were derived,
validated, and sealed before the holomorphic-twist literature for this problem
was consulted. The present section compares that sealed result with the
independent prediction of Budzik and Kulp~\cite{BK}, who compute the one-loop
correction to the supersymmetry algebra of the holomorphic twist of
$\mathcal N=4$ super-Yang--Mills theory. The comparison is an external
cross-check: no coefficient of~\cite{BK} enters the derivation of
Section~\ref{sec:superspace}, and everything concluded here is collected in
the admissibility statement of Subsection~\ref{subsec:ht:verdict}.

\subsection{The external target}\label{subsec:ht:target}

In the holomorphic twist~\cite{BK,Costello} the $\mathcal N=1$ vector
multiplet becomes a holomorphic $bc$ system and each chiral multiplet a
holomorphic $\beta\gamma$ system on $\mathbb C^{2}$. The Batalin--Vilkovisky
fields are
\begin{equation}\label{eq:ht:fields}
\begin{aligned}
b &= b^{(0)}+b^{(1)}+b^{(2)} \in
\Omega^{0,\bullet}(\mathbb C^{2},\mathfrak g^{\vee}) \ , &
c &= c^{(0)}+c^{(1)}+c^{(2)} \in
\Omega^{0,\bullet}(\mathbb C^{2},\mathfrak g)[1] \ , \\
\beta_{I} &= \beta_{I}^{(0)}+\beta_{I}^{(1)}+\beta_{I}^{(2)} \in
\Omega^{0,\bullet}(\mathbb C^{2},R^{\vee})[1] \ , &
\gamma^{I} &= \gamma^{I(0)}+\gamma^{I(1)}+\gamma^{I(2)} \in
\Omega^{0,\bullet}(\mathbb C^{2},R) \ ,
\end{aligned}
\end{equation}
with action (their eq.~(3.1))
\begin{equation}\label{eq:ht:action}
S_{\mathrm{BV}}=\int_{\mathbb C^{2}}d^{2}z\,
\Big\{\Tr\Big[b\Big(\bar\partial c+\frac12[c,c]\Big)\Big]
+\beta_{I}\big(\bar\partial\gamma^{I}+c\,\gamma^{I}\big)
+W[\gamma]\Big\} \ .
\end{equation}
The tree differential acts as
\begin{equation}\label{eq:ht:Q0fields}
\begin{aligned}
Q_{0}\,c &= \frac12\,[c,c] \ , &
Q_{0}\,b &= [c,b]-\mu(\beta,\gamma) \ , \\
Q_{0}\,\gamma^{I} &= c\,\gamma^{I} \ , &
Q_{0}\,\beta_{I} &= \beta_{I}\,c+\partial_{\gamma^{I}}W \ ,
\end{aligned}
\end{equation}
with $\mu(\beta,\gamma)$ the moment map. The letter gradings are
$|b|=|\gamma|=0$ and $|c|=|\beta|=1$ mod $2$, and normal ordering is
supercommutative on letters. The untwisted identifications are
$\gamma\sim\bar\phi$, $\beta\sim\psi_{+}$,
$\partial_{\dot\alpha}c\sim\bar\lambda_{\dot\alpha}$, $b\sim F_{++}$, and
$\partial_{\dot\alpha}\sim D_{+\dot\alpha}$, consistent with the component
table of Section~\ref{sec:component}. The Lie algebra is normalized as (their
eq.~(3.11))
\begin{equation}\label{eq:ht:lie}
[t_{A},t_{B}]=f_{AB}{}^{C}\,t_{C} \ ,\qquad
\Tr(t_{A}t_{B})=\kappa_{H}\,\delta_{AB} \ ,\qquad
\Tr(t_{A}[t_{B},t_{C}])=\kappa_{H}\,f_{ABC} \ ,
\end{equation}
with $\kappa_{H}$ the dual Coxeter number; the missing factor of $i$ relative
to our convention $[T_{A},T_{B}]=if_{AB}{}^{C}T_{C}$ is supplied by the frame
map \eqref{eq:ht:frames} below. For the $\mathcal N=4$ theory the matter
representation is $R=\mathfrak g\otimes\mathbb C^{3}$, the superpotential is
$W[\gamma]=\Tr\,\gamma^{1}[\gamma^{2},\gamma^{3}]
=\frac{1}{3!}\,\varepsilon_{IJK}\kappa_{H}f_{ABC}\,
\gamma^{IA}\gamma^{JB}\gamma^{KC}$, the twisting supercharge is $Q^{4}_{-}$,
and $\gamma^{I}\sim\Phi^{4I}$, $\beta_{I}\sim\Psi_{I+}$,
$\partial_{\dot\alpha}c\sim\bar\Psi^{4}_{\dot\alpha}$, $b\sim F_{++}$. The
$1/16$-BPS fields assemble into a single superfield on $\mathbb C^{2|3}$
(their eq.~(3.56)),
\begin{equation}\label{eq:ht:C}
C=c+\theta_{I}\gamma^{I}
+\frac12\,\varepsilon^{IJK}\theta_{I}\theta_{J}\beta_{K}
+\theta_{1}\theta_{2}\theta_{3}\,b \ ,
\end{equation}
in terms of which the action is a holomorphic Chern--Simons theory,
\begin{equation}\label{eq:ht:hcs}
S_{\mathrm{BV}}=\frac12\int_{\mathbb C^{2|3}}d^{2}z\,d^{3}\theta\,
\Tr\Big[C\Big(\bar\partial C+\frac13[C,C]\Big)\Big] \ .
\end{equation}
The classical supercharge obeys the tree identity (their eq.~(1.9))
\begin{equation}\label{eq:ht:Q0C}
Q_{0}\,C(\theta)=\frac12\,[C(\theta),C(\theta)] \ .
\end{equation}
Their one-loop result is a compact prediction for the loop-corrected
supercharge $Q_{1}$ (their eqs.~(1.10) and~(3.69)),
\begin{equation}\label{eq:ht:Q1C}
Q_{1}\big(C^{A}(\theta)\,C^{B}(\theta')\big)
=-\kappa_{H}^{2}\,f_{ACD}f_{BCE}\,
(\theta_{1}-\theta_{1}')(\theta_{2}-\theta_{2}')(\theta_{3}-\theta_{3}')\,
\partial_{\dot\alpha}C^{D}(\theta)\,
\partial^{\dot\alpha}C^{E}(\theta') \ .
\end{equation}
(The introduction of~\cite{BK} prints the coefficient $-1/4$ where its main
text prints $-\kappa_{H}^{2}$; this is the second source inconsistency, see
Subsection~\ref{subsec:ht:conflicts}.) Expanded in the components of
\eqref{eq:ht:C}, the prediction reads (their eqs.~(3.63)--(3.68))
\begin{equation}\label{eq:ht:componentpairs}
\begin{aligned}
Q_{1}(b^{A}c^{B})
&=\kappa_{H}^{2}f_{ACD}f_{BCE}\,
\partial_{\dot\alpha}c^{D}\partial^{\dot\alpha}c^{E} \ ,\\
Q_{1}(b^{A}b^{B})
&=\kappa_{H}^{2}f_{ACD}f_{BCE}
\big[\partial_{\dot\alpha}c^{D}\partial^{\dot\alpha}b^{E}
-\partial_{\dot\alpha}b^{D}\partial^{\dot\alpha}c^{E}
+\partial_{\dot\alpha}\beta_{I}^{D}\partial^{\dot\alpha}\gamma^{IE}
-\partial_{\dot\alpha}\gamma^{ID}\partial^{\dot\alpha}\beta_{I}^{E}\big] \ ,\\
Q_{1}(\beta_{I}^{A}\gamma^{JB})
&=\kappa_{H}^{2}\delta_{I}{}^{J}f_{ACD}f_{BCE}\,
\partial_{\dot\alpha}c^{D}\partial^{\dot\alpha}c^{E} \ ,\\
Q_{1}(\beta_{I}^{A}\beta_{J}^{B})
&=\kappa_{H}^{2}\varepsilon_{IJK}f_{ACD}f_{BCE}
\big[\partial_{\dot\alpha}c^{D}\partial^{\dot\alpha}\gamma^{KE}
-\partial_{\dot\alpha}\gamma^{KD}\partial^{\dot\alpha}c^{E}\big] \ ,\\
Q_{1}(b^{A}\gamma^{IB})
&=\kappa_{H}^{2}f_{ACD}f_{BCE}
\big[\partial_{\dot\alpha}c^{D}\partial^{\dot\alpha}\gamma^{IE}
-\partial_{\dot\alpha}\gamma^{ID}\partial^{\dot\alpha}c^{E}\big] \ ,\\
Q_{1}(\beta_{I}^{A}b^{B})
&=\kappa_{H}^{2}f_{ACD}f_{BCE}
\big[\partial_{\dot\alpha}\beta_{I}^{D}\partial^{\dot\alpha}c^{E}
+\partial_{\dot\alpha}c^{D}\partial^{\dot\alpha}\beta_{I}^{E}
+\varepsilon_{IJK}\partial_{\dot\alpha}\gamma^{JD}
\partial^{\dot\alpha}\gamma^{KE}\big] \ ,
\end{aligned}
\end{equation}
while $Q_{1}(cc)=Q_{1}(\gamma\gamma)=Q_{1}(\gamma c)=Q_{1}(\beta c)=0$.
In the $L_{\infty}$ language of~\cite{BK} the $(\ell+2)$-ary bracket is the
$\ell$-loop correction, with $\{\mathcal I,\cdot\}_{0}=Q_{0}$ and
$\frac{1}{2!}\{\mathcal I,\mathcal I,\cdot\}_{0}=Q_{1}$; the Maurer--Cartan
identities give $\{Q_{0},Q_{1}\}=0$ and $Q_{1}^{2}+\{Q_{0},Q_{2}\}=0$. Hence
$Q_{1}$ is nilpotent only on $Q_{0}$-cohomology and higher corrections
$Q_{n}$ are necessarily present: the construction of~\cite{BK} is not
one-loop exact, and the comparison below concerns its one-loop operation
$Q_{1}$.

\subsection{The dictionary}\label{subsec:ht:dictionary}

Let $R_{A}{}^{P}$ be a color frame, $U_{I}{}^{r}$ a flavor frame,
$S_{\dot a}{}^{\dot b}$ a dotted-spin frame, and $\rho\neq0$ an overall
scale. The exact map between the letters
$b,\beta_{r},\gamma^{r},\partial_{\dot\alpha}c$ of
Section~\ref{sec:superspace} and the alphabet of~\cite{BK} is
\begin{equation}\label{eq:ht:dictionary}
\begin{aligned}
b^{A}&\longmapsto\rho\,R_{A}{}^{P}\,b^{P} \ ,\\
\beta_{I}^{A}&\longmapsto\rho\,R_{A}{}^{P}U_{I}{}^{r}\,
\beta_{r}^{P} \ ,\\
\gamma^{IA}&\longmapsto\rho\,R_{A}{}^{P}(U^{-1})_{r}{}^{I}\,
\gamma^{rP} \ ,\\
\partial_{\dot a}c^{A}&\longmapsto\rho\,R_{A}{}^{P}
S_{\dot a}{}^{\dot b}\,\partial_{\dot b}c^{P} \ ,
\end{aligned}
\end{equation}
where the left-hand sides are the fields of \eqref{eq:ht:fields} and the
right-hand sides use the letters of Section~\ref{sec:superspace}. Every row
carries the common scale $\rho$: the factors $-i/\sqrt2$ and $1/\sqrt2$
familiar from the raw superfield composites are precisely those absorbed
into the letter definitions of Definition~\ref{def:ss:letters}, so that in
the letters of this paper the map is uniform. The $\gamma$ row carries the
inverse flavor frame and the $\partial c$ row the spin frame. At $\rho=1$
with identity frames the map reduces to the identity,
$b\mapsto b$, $\beta_{I}\mapsto\beta_{r}\delta_{I}{}^{r}$,
$\gamma^{I}\mapsto\gamma^{r}\delta^{I}{}_{r}$, and
$\partial_{\dot a}c\mapsto\partial_{\dot a}c$. The frames enter the color
and trace relations through
\begin{equation}\label{eq:ht:frames}
R_{A}{}^{P}R_{B}{}^{Q}\,\delta_{PQ}=\delta_{AB} \ ,\qquad
t_{A}=-i\,R_{A}{}^{P}T_{P} \ ,\qquad
f_{AB}{}^{C}=R_{A}{}^{P}R_{B}{}^{Q}f_{PQ}{}^{R}(R^{-1})_{R}{}^{C} \ ,\qquad
\Tr_{H}=-\kappa_{H}\Tr \ ,
\end{equation}
where $\Tr_{H}$ is the trace of \eqref{eq:ht:lie} and $\Tr$ ours, with
$\kappa_{AB}=\delta_{AB}$. Applied to \eqref{eq:ht:C}, the map sends
\begin{equation}\label{eq:ht:Cimage}
C(\theta)\;\longmapsto\;
\rho\Big[c+\theta_{r}\gamma^{r}
+\frac{1}{2}\,\varepsilon_{rst}\theta_{r}\theta_{s}\beta_{t}
+\theta_{1}\theta_{2}\theta_{3}\,b\Big] \ ;
\end{equation}
the first entry is the ghost (the formal ghost placeholder of the letter
algebra), which takes part in the twisted theory but not in the physical
letter pairs compared below. Matching the one-loop coefficients of the two
formalisms on the gaugino ($\cW_{+}$) component gives the loop-scale relation
\begin{equation}\label{eq:ht:loopscale}
\h_{\mathrm{HT}}\,\kappa_{H}^{2}
=\frac{\h\,g^{2}}{32\sqrt2\,\pi^{2}}
=\frac{1}{2\sqrt2}\,\frac{\h g^{2}}{16\pi^{2}} \ .
\end{equation}

\begin{remark}[conditional status of \eqref{eq:ht:loopscale}]\label{rem:ht:conditional}
The relation \eqref{eq:ht:loopscale} is obtained under a conditional
total-differential scale in the dictionary analysis; it is a conditional
dictionary entry, not a theorem proved in this paper. The four physical
letter rows of \eqref{eq:ht:dictionary} are independent of this
conditionality, and the comparison of Subsection~\ref{subsec:ht:results} is
performed on output words only. A single $\kappa_{H}$ presupposes a simple
gauge algebra; for a general reductive gauge algebra the trace relations hold
factorwise.
\end{remark}

\subsection{Comparison results}\label{subsec:ht:results}

The comparison proceeds word by word. Each ordered pair of the nine letters
$(b,\beta_{r},\gamma^{r},\partial_{\dot a}c)$ is mapped by
\eqref{eq:ht:dictionary}; the predicted output monomials are translated into
the letters of Section~\ref{sec:superspace}; and the result is compared with
the sealed coefficient table over $\mathbb Q(i,\sqrt2)$. With the
identification $Q_{0}\leftrightarrow-\tfrac12\nabla_{-}$ of
Section~\ref{sec:superspace}, the quantity compared with $Q_{1}(XY)$ is the
one-loop anomalous variation $\nabla_{-}(XY)|_{\text{1-loop}}$, normalized by
\eqref{eq:ht:loopscale}.

\begin{proposition}[ordered-word comparison]\label{thm:ht:81}
For every ordered pair of letters, the sealed coefficient word of
Section~\ref{sec:superspace}, its independent expansion, and the translated
holomorphic-twist word agree exactly:
\begin{equation}\label{eq:ht:81}
N_{\mathrm{pairs}}=81 \ ,\qquad
N_{\mathrm{agree}}=81 \ ,\qquad
N_{\mathrm{disagree}}=0 \ ,\qquad
81=29_{\mathrm{nonzero}}+52_{\mathrm{zero}} \ .
\end{equation}
The nonzero census decomposes as
\begin{equation}\label{eq:ht:29}
29=\underset{(b,b)}{1}
+\underset{(b,\beta_{r}),(\beta_{r},b)}{6}
+\underset{(b,\gamma^{r}),(\gamma^{r},b)}{6}
+\underset{(\beta_{r},\beta_{s}),\,r\neq s}{6}
+\underset{(\beta_{r},\gamma^{r}),(\gamma^{r},\beta_{r})}{6}
+\underset{(b,\partial_{\dot a}c),(\partial_{\dot a}c,b)}{4} \ ,
\end{equation}
comprising $8+6\cdot4+6\cdot2+6\cdot2+6\cdot1+4\cdot2=70$ nonzero output
words; the remaining $52$ ordered pairs vanish on both sides.
\end{proposition}

The mixed rows $(b,\partial_{\dot a}c)$ and $(\partial_{\dot a}c,b)$ display
the intrinsic split $K_{(e_{\dot a});0}=\tfrac23$ and
$K_{(e_{\dot a});(e_{\dot a})}=\tfrac13$ for a unit dotted derivative
$e_{\dot a}$, the two orderings being exchanged by the Leibniz rule; this is
the $\tfrac13$--$\tfrac23$ distribution already recorded in the channel table
of Section~\ref{sec:superspace}. At arbitrary derivative order the comparison
uses the Appendix-B tower of~\cite{BK} (their eq.~(B.5)),
\begin{equation}\label{eq:ht:Cmn}
m!\,n!\,C_{mn}
=\frac{1}{m+n+2}\sum_{k=0}^{m}\sum_{\ell=0}^{n}
\frac{1}{k+\ell+1}\binom{m}{k}\binom{n}{\ell}\,
(p_{1})_{1}^{k}(p_{1})_{2}^{\ell}\,
(p_{2})_{1}^{m-k}(p_{2})_{2}^{n-\ell} \ ,
\end{equation}
whose component formulas (their eqs.~(B.6)--(B.11)) multiply this tower by
$\kappa_{H}^{2}f_{ACD}f_{BCE}$ with the same flavor structures as the
zero-derivative pairs \eqref{eq:ht:componentpairs}; we write their holomorphic
form variables as $p_{i}$, $p_{1}+p_{2}+p_{3}=0$, to avoid a clash of
notation.

\begin{proposition}[all-orders kernel identity]\label{thm:ht:kernel}
For all $m,n\in\mathbb Z_{\ge0}$ and $0\le k\le m$, $0\le\ell\le n$,
\begin{equation}\label{eq:ht:kernel}
K_{m,n;k,\ell}
=\frac{2\,\binom{m}{k}\binom{n}{\ell}}{(m+n+2)\,(k+\ell+1)}
=2\,T^{\mathrm{HT,printed}}_{m,n;k,\ell}
=T^{\mathrm{HT,corrected}}_{m,n;k,\ell} \ ,
\end{equation}
where $K_{m,n;k,\ell}$ is the derivative kernel of the one-loop variation of
Section~\ref{sec:superspace},
\begin{equation}\label{eq:ht:Tprinted}
T^{\mathrm{HT,printed}}_{m,n;k,\ell}
=\frac{\binom{m}{k}\binom{n}{\ell}}{(m+n+2)\,(k+\ell+1)}
\end{equation}
is the coefficient tower of \eqref{eq:ht:Cmn}, and
$T^{\mathrm{HT,corrected}}:=2\,T^{\mathrm{HT,printed}}$. At the origin
$T^{\mathrm{HT,printed}}_{0,0;0,0}=\tfrac12$ while $K_{0,0;0,0}=1$.
\end{proposition}

The factor $2$ in \eqref{eq:ht:kernel} is the $\Gamma(3)=2$ of the triangle
Feynman parametrization,
\begin{equation}\label{eq:ht:feynman}
\frac{1}{D_{0}D_{1}D_{2}}
=2\int_{\Delta_{2}}\frac{1}{(r^{2}+\Delta)^{3}} \ ,
\qquad 2\int_{\Delta_{2}}1=1 \ ,
\end{equation}
not an orientation multiplicity. The triangle operator of~\cite{BK} (their
eq.~(3.8)), with $\partial=(\partial_{1},\partial_{2})$ the holomorphic
derivative on $\mathbb C^{2}$,
\begin{equation}\label{eq:ht:dtri}
\mathcal D^{\triangle}_{w,z}(f,g)
:=\sum_{m,n=0}^{\infty}\frac{1}{(n+m+2)!}\,
e^{z\cdot\partial}\big((w-z)\cdot\partial\big)^{n}\,
\partial_{\dot\alpha}f(y)\,
e^{w\cdot\partial}\big((z-w)\cdot\partial\big)^{m}\,
\partial^{\dot\alpha}g(y)\,\Big|_{y=0} \ ,
\end{equation}
generates the generic shifted pairs, each $Q_{1}$ pair action being
$\kappa_{H}^{2}f_{ACD}f_{BCE}$ times the corresponding
$\mathcal D^{\triangle}$ bilinear; its printed zero-shift value (their
eq.~(3.9)),
\begin{equation}\label{eq:ht:dtri00}
\mathcal D^{\triangle}_{0,0}(f,g)
=\frac12\,\partial_{\dot\alpha}f\,\partial^{\dot\alpha}g \ ,
\end{equation}
misses precisely the prefactor $2$ of \eqref{eq:ht:feynman}, while the
printed component formulas \eqref{eq:ht:componentpairs} carry the
complete-amplitude normalization (see
Subsection~\ref{subsec:ht:conflicts}). The identity \eqref{eq:ht:kernel} is
proved symbolically for all $m,n\ge0$ and checked coefficient by coefficient
on the finite rectangle $0\le m,n\le8$:
\begin{equation}\label{eq:ht:2025}
\sum_{m=0}^{8}\sum_{n=0}^{8}(m+1)(n+1)
=\Big(\sum_{j=1}^{9}j\Big)^{2}=45^{2}=2025
\quad\text{kernel checks, all passing }(2025/2025)\ .
\end{equation}
Multiplication by the $70$ nonzero base output words lifts this to the full
derivative tower,
\begin{equation}\label{eq:ht:141750}
2025\times70=141750
\quad\text{lifted coefficient checks, all passing }(141750/141750)\ ,
\end{equation}
while the $52$ vanishing rows remain zero at every derivative order. The
independent physical verification suite accepts the physical one-loop anomaly
sector with $33/33$ checks passed and $0$ failed.

\subsection{Two source-internal inconsistencies and their adjudication}\label{subsec:ht:conflicts}

The source~\cite{BK} contains two internal normalization conflicts. Both are
quarantined as matters of source provenance: the coefficients of this paper
are derived independently in Section~\ref{sec:superspace}, and the comparison
above uses only output words, so no contested normalization of~\cite{BK} is
imported.

\paragraph{The factor two.} Printed inside~\cite{BK} are: the zero-shift
kernel \eqref{eq:ht:dtri00}, carrying $1/2=1/\Gamma(3)$; the shifted-pair
formulas multiplying $\mathcal D^{\triangle}$ by $\kappa_{H}^{2}$, hence
$\kappa_{H}^{2}/2$ at $z=w=0$; the zero-shift component formulas
\eqref{eq:ht:componentpairs}, printing $\kappa_{H}^{2}$ times the derivative
bilinears; and the tower \eqref{eq:ht:Cmn}, printing
$\kappa_{H}^{2}/(m+n+2)=\kappa_{H}^{2}/2$ at $m=n=0$. The evaluated branches
give
\begin{equation}\label{eq:ht:branches}
\mathcal I_{\triangle}[p;0]
=\mathcal D^{\triangle}_{0,0}
=C_{00}
=\frac12 \ ,
\qquad
\frac{c_{\mathrm{component}}}{c_{\mathrm{triangle}}}
=\frac{c_{\mathrm{compact}}}{c_{\mathrm{triangle}}}
=2 \ ,
\end{equation}
with $\mathcal I_{\triangle}[p;0]=\tfrac12\,p_{1}\wedge p_{2}$ the zero-shift
master triangle (their eqs.~(2.35) and~(2.39)). The adjudication selects the
complete-amplitude branch: the zero-shift component block and the compact
block of~\cite{BK} carry the complete normalization, while its
master-integral, generic shifted-pair, and Appendix-B branches are uniformly
smaller by the factor $2$ and must be multiplied by $2$ before comparison
with the complete amplitude, consistently with \eqref{eq:ht:kernel} holding
at every derivative degree. The conflict remains unresolved inside the
source: the classification recorded here repairs nothing in~\cite{BK} and
changes no relative coefficient of Section~\ref{sec:superspace}.

\paragraph{The compact coefficient.} The introduction of~\cite{BK} prints
the compact prediction \eqref{eq:ht:Q1C} with coefficient $-1/4$; its main
text prints $-\kappa_{H}^{2}$. The two displays agree only if
\begin{equation}\label{eq:ht:conflict2}
\kappa_{H}^{2}=\frac14 \ .
\end{equation}
The independently fixed coefficient of Section~\ref{sec:superspace} selects
the main-text branch $-\kappa_{H}^{2}$ and rejects the introductory display
unless $\kappa_{H}^{2}=1/4$; again this classifies the conflict without
repairing the source.

\subsection{The admissibility verdict}\label{subsec:ht:verdict}

\begin{theorem}[admissibility of the holomorphic-twist comparison]\label{thm:ht:verdict}
Under the dictionary \eqref{eq:ht:dictionary}, the identification
$Q_{0}\leftrightarrow-\tfrac12\nabla_{-}$ of
Section~\ref{sec:superspace}, and the loop-scale relation
\eqref{eq:ht:loopscale}, the direct superspace calculation of
Section~\ref{sec:superspace} exactly reproduces the holomorphic-twist
prediction for the one-loop semi-chiral anomaly: the $81$ ordered output
words agree as in \eqref{eq:ht:81} ($29$ nonzero, $52$ zero), and the
derivative kernels agree to all orders,
\begin{equation}\label{eq:ht:verdict}
K_{m,n;k,\ell}=T^{\mathrm{HT,corrected}}_{m,n;k,\ell}
\qquad\text{for all }m,n\ge0 \ ,
\end{equation}
i.e.\ the sealed one-loop coefficient table coincides with the
factor-two-corrected one-loop supercharge $Q_{1}$ of~\cite{BK} on every
output word, with no mismatch.
\end{theorem}

\begin{remark}[agreement of totals is not an absence proof]\label{rem:ht:cr23}
Theorem~\ref{thm:ht:verdict} compares total coefficients. Writing
$C_{\mathrm{complete}}=C_{\mathrm{tabulated}}+\sum_{j\in\mathcal O}c_{j}$ for
the contributions $c_{j}$ of any set $\mathcal O$ of omitted Wick routings,
\begin{equation}\label{eq:ht:cr23}
C_{\mathrm{tabulated}}=C_{\mathrm{HT}}
\quad\text{does not imply}\quad
\mathcal O=\varnothing
\ \text{or}\ c_{j}=0\ \text{for every }j \ :
\end{equation}
cancellations such as $c_{1}=t$, $c_{2}=-t$ with $t\neq0$ leave the total
unchanged. Equality of the total coefficient with the holomorphic-twist
prediction therefore does not prove termwise absence of omitted Wick
routings; the agreement is a final cross-check, not a target-blind absence
proof.
\end{remark}

\begin{remark}[the adversarial census remains open]\label{rem:ht:open}
The completeness question for the graph census is not settled by
Theorem~\ref{thm:ht:verdict} and carries the status
\texttt{OPEN\_FINAL\_CENSUS\_17\_CANDIDATES}: the
ghost--Nielsen--Kallosh sector has not run, the zero/cut-orbit sector has not
run, the automated three-lens refutation has not run, and four
physics-relevant Wick routings remain deferred to a future per-routing
derivation from the locked action, without reading the target.
\end{remark}

\begin{remark}[the accepted coefficients are unchanged]\label{rem:ht:sealed}
The accepted coefficient table of Section~\ref{sec:superspace} rests on the
target-blind Schwinger-cut derivation presented there; it was read,
validated, and sealed (SHA-256 \texttt{ec77327d\ldots b3016d}) before the
source of~\cite{BK} was consulted, and it is unaffected by the open census:
within the physical one-loop anomaly sector no coefficient row is
unresolved.
\end{remark}
